% ============================================================
% ［架空の記入例・提出不可］対象データ，判定条件および評価設計はすべて架空である．
% 通常のchapters/03-problem.texには自動で読み込まれない．
% ============================================================
\chapter{問題設定}

\begin{center}
  \UsamiFictionalExampleNotice
\end{center}

\section{対象と仮定}

入力は，固定カメラで撮影した低照度の単眼RGB動画とする．各フレームに既存の2次元姿勢推定器を適用し，検出人物ごとに$J$個の関節座標と信頼度を得る．カメラ校正値，深度画像および時間的に前後するフレームは，骨格の採否判定には用いない．姿勢推定器の学習および人物領域の検出は本研究の対象外とし，少なくとも胴体長を定義する関節が出力されることを仮定する．

低照度条件では人物そのものが未検出となる場合もあるが，本研究が判定できるのは推定器から出力された骨格だけである．したがって，未検出人物の回復と，出力済み骨格に含まれる誤推定の除外を区別する．また，評価時に正解関節が与えられるが，品質スコアの算出時には正解を使用しない．

\section{記号と定式化}

フレーム$t$で検出された人物$n$の関節$j$について，推定座標を$\hat{\bm{p}}_{tnj}\in\mathbb{R}^{2}$，信頼度を$c_{tnj}\in[0,1]$とする．骨格辺集合を$\mathcal{E}$とし，胴体長で正規化した辺$(j,k)\in\mathcal{E}$の長さを$\ell_{tnjk}$とする．学習系列から求めた辺長の基準値を$\mu_{jk}$とし，評価系列の情報は基準値の算出に用いない．

骨格選別器$g$は，関節座標と信頼度を入力として採用フラグ$a_{tn}\in\{0,1\}$を返す．

\begin{equation}
  a_{tn}
  =g\!\left(
    \{\hat{\bm{p}}_{tnj},c_{tnj}\}_{j=1}^{J},
    \{\mu_{jk}\}_{(j,k)\in\mathcal{E}}
  \right).
  \label{eq:problem-selection}
\end{equation}

式\eqref{eq:problem-selection}において，$a_{tn}=1$なら骨格を後段処理へ渡し，$a_{tn}=0$なら誤推定候補として除外する．目標は，正しい骨格をできる限り残しながら，大きな関節位置誤差または人体構造からの逸脱をもつ骨格を除外することである．第4章では，$g$を関節信頼度と正規化辺長の逸脱から構成する．

\section{評価課題}

正解関節座標を$\bm{p}_{tnj}$，人物の基準長を$s_{tn}$とする．関節$j$が正しい条件を

\begin{equation}
  \left\lVert
    \hat{\bm{p}}_{tnj}-\bm{p}_{tnj}
  \right\rVert_{2}
  \leq 0.2s_{tn}
  \label{eq:problem-pck}
\end{equation}

とし，採用骨格に含まれる可視関節について式\eqref{eq:problem-pck}を満たす割合をPCK@0.2として集計する．同時に，全検出骨格数に対する$a_{tn}=1$の骨格数の割合を採用率とする．前者は残した骨格の位置精度，後者は除外による情報欠損を表すため，どちらか一方だけでは選別器を評価しない．

判定閾値は検証系列で固定し，通常照明，低照度および動きぼけを含む評価系列へ同じ値を適用する．主要比較は，除外なし，平均関節信頼度による除外，構造的一貫性を統合した提案法の3条件とする．本章の入力条件，記号の運用，閾値決定および評価課題は，すべて架空の記入例である．
